\documentclass[11pt]{article}

\usepackage[margin=1in]{geometry}
\usepackage{booktabs}
\usepackage{hyperref}
\usepackage{xcolor}
\usepackage{listings}
\emergencystretch=3em

\hypersetup{
  colorlinks=true,
  linkcolor=blue,
  urlcolor=blue
}

\lstdefinestyle{code}{
  basicstyle=\ttfamily\small,
  breaklines=true,
  frame=single,
  columns=fullflexible,
  keepspaces=true
}
\lstset{style=code}

\title{Progress Report: Agentic Python Verifier}
\author{Harry Wang \and Eric Shang Nan Chen}
\date{May 6, 2026}

\begin{document}

\maketitle

\section{Project Overview}

The project is a proof-driven code generation pipeline for Python. Instead of
generating Python first and checking it afterward, the system first asks a
language-model synthesis agent to produce a PlusCal algorithm, a candidate
inductive invariant, and a target safety property from a natural-language
requirement. The formal verifier then translates the PlusCal module to TLA+,
checks three proof obligations with TLC, and returns structured
counterexamples to the synthesis agent when a proof fails. Only after all
three obligations pass does a refinement agent translate the verified PlusCal
artifact into executable Python with runtime assertions derived from the
verified invariant.

The current repository contains a working first pipeline implementation,
command-line interface, formal checking layer, language-model tool contracts,
offline tests, integration-test scaffolding, and two example benchmark
requirements.

\section{Implemented Pipeline}

The implemented pipeline is centered in \texttt{src/main.py}. It accepts a
\texttt{TaskRequest}, loads runtime settings, creates output directories, and
then runs the following loop:

\begin{enumerate}
  \item The synthesis agent proposes a tuple
  \texttt{(PlusCal, Inv, Property, finite constants)}.
  \item The verifier checks initiation, consecution, and property implication.
  \item If any obligation fails, the failing TLC result is serialized and sent
  back to the repair agent through Anthropic tool-result history.
  \item If all obligations pass, the refinement agent emits a Python module.
  \item Verified artifacts are written to \texttt{generated/tla} and
  \texttt{generated/python}, with per-iteration scratch files in
  \texttt{generated/work}.
\end{enumerate}

The pipeline returns a structured \texttt{PipelineResult} with either
\texttt{verified} or \texttt{unverified} status. Verification failures do not
silently pass through to Python generation; Python is emitted only when the
full proof bundle passes.

\section{Formal Verification Progress}

The verifier in \texttt{src/agents/verifier.py} implements three model-checking
obligations for a proposal with initial predicate \texttt{Init}, transition
relation \texttt{Next}, inductive invariant \texttt{Inv}, and target property
\texttt{Property}.

\begin{table}[h]
\centering
\begin{tabular}{p{0.18\linewidth}p{0.24\linewidth}p{0.48\linewidth}}
\toprule
Obligation & Mathematical Form & TLC Encoding \\
\midrule
Initiation & \texttt{Init => Inv} &
\texttt{SPECIFICATION Spec}, \texttt{INVARIANT Inv} \\
Consecution & \texttt{Inv /\ Next => Inv'} &
Auxiliary \texttt{IInit == Inv}, \texttt{ISpec == IInit /\ [][Next]\_vars} \\
Property implication & \texttt{Inv => Property} &
Auxiliary \texttt{PInit == Inv}, unchanged-step \texttt{PSpec} \\
\bottomrule
\end{tabular}
\caption{Implemented proof obligations.}
\end{table}

The \texttt{src/formal/tlc\_config.py} module generates the TLC configuration
files and the two auxiliary TLA+ modules used for consecution and property
checking. The property check enumerates all invariant states by using
\texttt{PInit == Inv} and an unchanged transition relation, then asks TLC to
check \texttt{Property} as an invariant over those states.

The implementation also wraps the two required TLA+ tools:

\begin{itemize}
  \item \texttt{src/formal/pcal\_translator.py} invokes
  \texttt{pcal.trans} from \texttt{tla2tools.jar} and checks that a translation
  block was inserted.
  \item \texttt{src/formal/tla\_runner.py} invokes \texttt{tlc2.TLC} with a
  configurable timeout, automatic workers, and deadlock checking enabled.
\end{itemize}

\section{Counterexample Feedback}

Structured counterexample feedback is implemented in
\texttt{src/formal/counterexample\_parser.py}. The parser classifies TLC runs
as \texttt{passed}, \texttt{failed}, \texttt{timeout}, or \texttt{error}. For
invariant and initial-state violations it extracts:

\begin{itemize}
  \item the failing obligation, such as \texttt{consec};
  \item the violated predicate, such as \texttt{Inv};
  \item each trace state index and action;
  \item variable assignments from TLC trace lines.
\end{itemize}

This information is then serialized by \texttt{SynthesisAgent.repair}. For
example, a consecution failure may be sent back to the model as:

\begin{lstlisting}
- consec: FAILED
    violated_predicate: Inv
    State 1 <Initial predicate>: counter=11
\end{lstlisting}

This is a major milestone because the system is not limited to one-shot
generation. It has a repair loop with formal counterexamples as the feedback
signal.

\section{Language-Model Agent Progress}

The language-model boundary is implemented with explicit Anthropic tool
schemas in \texttt{src/llm/tools.py}. Three tools are currently defined:

\begin{itemize}
  \item \texttt{propose\_pluscal\_with\_invariant}: emits the initial PlusCal
  module, invariant, property, and finite constant domains.
  \item \texttt{repair\_after\_counterexample}: emits a revised proposal after
  a failed proof obligation.
  \item \texttt{emit\_python\_module}: emits Python source and an assertion map
  from invariant clauses to Python checks.
\end{itemize}

The prompt layer in \texttt{src/llm/prompts.py} gives the synthesis model
specific formal-methods constraints. In particular, it requires small finite
constant domains, explicit variable membership in \texttt{Inv}, and explicit
enumeration of the PlusCal \texttt{pc} variable. These details are important
because the consecution and property checks use \texttt{Inv} as an initial
state predicate, so TLC must be able to enumerate states from the invariant
alone.

The Anthropic client wrapper in \texttt{src/llm/anthropic\_client.py} supports
prompt caching through ephemeral cache-control system blocks and retries once
on overloaded-model errors using a fallback model.

\section{Python Refinement Progress}

The refinement stage is implemented in \texttt{src/agents/refine\_agent.py}.
It extracts top-level conjuncts from the verified invariant and passes them to
the refinement model. The emitted Python module is required to include runtime
assertions at the entry and exit of state-mutating functions and to provide an
\texttt{assertion\_map} for traceability.

For example, given this invariant:

\begin{lstlisting}
Inv == counter >= 0
    /\ counter <= 10
    /\ counter \in 0..10
\end{lstlisting}

the refinement prompt sends the model the individual conjuncts:

\begin{lstlisting}
- counter >= 0
- counter <= 10
- counter \in 0..10
\end{lstlisting}

The intended Python result is an implementation that carries corresponding
assertions, such as:

\begin{lstlisting}[language=Python]
def increment(counter: int) -> int:
    assert counter >= 0
    assert counter <= 10
    if counter < 10:
        counter += 1
    assert counter >= 0
    assert counter <= 10
    return counter
\end{lstlisting}

\section{Examples}

Two benchmark prompts have been added in the \texttt{examples} directory:

\begin{lstlisting}
Implement a bounded counter from 0 to 10.
\end{lstlisting}

\begin{lstlisting}
Implement a bank transfer between two accounts that never allows a negative
balance and preserves the total sum of money.
\end{lstlisting}

The test fixtures include a good bounded-counter TLA+ module and a deliberately
bad invariant module. The good fixture uses:

\begin{lstlisting}
Inv == /\ counter \in 0..MaxValue
       /\ pc \in {"Increment", "Done"}

Property == counter <= MaxValue
\end{lstlisting}

This shows the expected invariant style: the counter is bounded by explicit set
membership and the PlusCal control state is enumerated. The bad fixture uses
\texttt{counter \textbackslash in 0..(MaxValue-1)}, which is too strong after
the final increment. The integration test expects this module to fail an
initiation or consecution obligation with a counterexample.

\section{Command-Line Interface}

A Typer-based CLI is implemented in \texttt{src/cli.py}. It exposes:

\begin{itemize}
  \item \texttt{--output} for generated artifacts;
  \item \texttt{--max-iterations} with bounds from 1 to 20;
  \item \texttt{--model} for overriding the Anthropic model;
  \item \texttt{--verbose} for debug logging.
\end{itemize}

The CLI uses distinct exit codes: \texttt{0} for verified output, \texttt{1}
for exhausted repair iterations, and \texttt{2} for missing runtime
configuration such as \texttt{ANTHROPIC\_API\_KEY} or \texttt{TLA2TOOLS\_JAR}.

\section{Testing Progress}

The repository includes focused unit tests and mocked pipeline tests:

\begin{itemize}
  \item model validation tests for task requests, synthesis proposals,
  constants, repair proposals, and proof bundles;
  \item configuration tests for missing API keys and TLA+ tools;
  \item CLI smoke tests for help text and fail-fast configuration errors;
  \item counterexample parser tests using captured TLC output fixtures;
  \item TLC config generation tests for initiation, consecution, and property
  configurations;
  \item mocked synthesis-agent tests for tool-call extraction and repair
  feedback serialization;
  \item mocked refinement-agent tests for invariant-conjunct extraction;
  \item mocked end-to-end pipeline tests covering first-try success, repair
  then success, and exhausted iterations;
  \item gated integration tests for running real TLC against known good and
  known bad bounded-counter fixtures.
\end{itemize}

The offline test suite is designed to run without an Anthropic API key or
\texttt{tla2tools.jar}; the external dependencies are mocked at the boundary.
The TLC integration tests are explicitly marked and skipped unless a valid
TLA+ tools jar is configured.

\section{Current Limitations}

The prototype currently targets finite-domain inductive checking with TLC.
Unbounded inductive proofs would require a theorem prover such as TLAPS, which
is outside the current scope. The invariant must be enumerable as an initial
state predicate, which means every state variable needs an explicit finite
membership constraint. The refinement step depends on the language model to
translate verified PlusCal into faithful Python, so the runtime assertions
provide an additional guard but do not replace independent Python-level tests.

There is also inherent nondeterminism in live LLM calls. The current design
mitigates this by making the verified TLA+ module and proof results the trusted
intermediate artifact, while the repository tests mock the Anthropic boundary
for repeatability.

\section{Next Steps}

The next development milestones are:

\begin{enumerate}
  \item Run live end-to-end examples with a real Anthropic API key and
  \texttt{tla2tools.jar}, then save representative generated artifacts.
  \item Add Python-level tests for refined modules so generated runtime
  assertions are exercised after formal verification succeeds.
  \item Expand the benchmark set beyond bounded counters and bank transfers.
  \item Improve invariant-conjunct extraction so it handles nested TLA+
  expressions more robustly than the current regular-expression split.
  \item Add reporting or logs that summarize each repair iteration and the
  obligation that drove it.
\end{enumerate}

\section{Summary}

Substantial progress has been made toward a complete proof-driven Python code
generation system. The project now has a working architecture for
PlusCal/invariant co-synthesis, three formal TLC proof obligations, structured
counterexample-driven repair, verified-artifact output, Python refinement with
invariant-derived assertions, a usable CLI, examples, and broad offline test
coverage. The remaining work is mainly in live evaluation, broader benchmarks,
and hardening the refinement and reporting layers.

\end{document}
